\chapter{Migraine}
\index{Migraine}
\label{sec:Migraine}

\section{Overview}
Migraine is a primary headache disorder affecting approximately 15\% of the global population (1 billion people), with a female-to-male ratio of 3:1. It is the second leading cause of years lived with disability worldwide. This neurological disorder affects the central or peripheral nervous system, leading to progressive or episodic symptoms. It significantly impacts quality of life and often requires multidisciplinary care. Neurological disorders affect the central or peripheral nervous system, leading to a wide range of symptoms affecting movement, sensation, cognition, and autonomic function. The nervous system's complexity means that even small disruptions can cause significant functional impairment. Many neurological conditions are chronic and progressive, requiring long-term management and rehabilitation. Advances in neuroimaging and molecular diagnostics have improved understanding and treatment of these conditions.
\section{Causes}
The underlying mechanisms involve complex interactions between genetic predisposition and environmental factors. Disruption of normal physiological processes leads to the characteristic features of the condition. Multiple contributing factors may act synergistically to trigger or worsen the condition.

\begin{itemize}[noitemsep]
  \item This condition happens because of cortical spreading depression (CSD)---a wave of neuronal depolarization followed by prolonged suppression of cortical activity---triggering the trigeminovascular system with release of calcitonin gene-related peptide (CGRP), causing vasodilation and neurogenic inflammation
  \item Genetic factors play a role, with familial hemiplegic migraine linked to mutations in CACNA1A, ATP1A2, and SCN1A.
\end{itemize}
\section{Risk Factors}
\begin{itemize}[noitemsep]
  \item Genetic predisposition and family history
  \item Advancing age
  \item Lifestyle factors (diet, exercise, smoking, alcohol)
  \item Environmental exposures
  \item Underlying medical conditions
  \item Medication use
\end{itemize}
\section{Signs and Symptoms}

\subsection{Common symptoms}
\begin{itemize}[noitemsep]
  \item You may experience recurrent, unilateral (60\%), pulsating headache of moderate to severe intensity, lasting 4--72 hours, associated with nausea, vomiting, photophobia, and phonophobia
  \item Migraine with aura involves transient focal neurological symptoms preceding or accompanying the headache
  \item Chronic migraine is defined by headache on $\ge$15 days per month for >3 months.
  \item Pain or discomfort in the affected area
  \end{itemize}

\subsection{Emergency warning signs}
\begin{itemize}[noitemsep]
  \item Severe or worsening symptoms of migraine require immediate medical evaluation
\end{itemize}
\section{Progression and Stages}
The condition may progress through different stages of severity over time. Early stages often involve mild or subtle symptoms that may be overlooked. Moderate stages involve more noticeable symptoms affecting daily function. Advanced stages may involve significant impairment and complications.
\section{Complications}
\begin{itemize}[noitemsep]
  \item Disease progression leading to worsening symptoms
  \item Secondary infections or injuries
  \item Side effects of treatment
  \item Reduced quality of life
  \item Emotional and psychological impact
\end{itemize}
\section{Diagnosis}
Doctors usually diagnose this based on your symptoms and a physical exam based on ICHD-3 criteria. Acute Treatment focuses on NSAIDs or triptans for mild to moderate attacks and triptans, lasmiditan, or gepants for moderate to severe attacks. Preventive therapy is indicated for $\ge$4 migraine days per month and includes beta-blockers, amitriptyline, topiramate, candesartan, onabotulinumtoxin A for chronic migraine, and CGRP monoclonal antibodies. Neurological diagnosis combines clinical history and examination findings with neuroimaging (MRI, CT), electrophysiological studies (EEG, EMG, nerve conduction), and laboratory testing of blood and cerebrospinal fluid.

\begin{itemize}[noitemsep]
  \item Comprehensive medical history and physical examination
  \item Laboratory tests to assess organ function
  \item Imaging studies to visualize affected structures
  \item Specialized testing depending on the affected system
  \item Consultation with relevant specialists
\end{itemize}
\section{Prevention}
\begin{itemize}[noitemsep]
  \item Maintain a healthy lifestyle with balanced diet and exercise
  \item Avoid known risk factors
  \item Attend regular medical check-ups for early detection
  \item Follow recommended screening guidelines
\end{itemize}
\section{Treatment Options}
Pharmacologic therapy targeting specific neurotransmitter systems or disease mechanisms Physical, occupational, and speech therapy for functional maintenance Surgical interventions when appropriate (deep brain stimulation, tumor resection) Cognitive rehabilitation and memory support Pain management for neuropathic pain Assistive devices and mobility aids Lifestyle modifications including exercise, nutrition, and sleep hygiene Support services and caregiver education Regular neurologic monitoring and treatment adjustment Clinical trials for emerging therapies

\begin{itemize}[noitemsep]
  \item Medications to manage symptoms and address underlying mechanisms
  \item Lifestyle modifications and self-care measures
  \item Physical therapy and rehabilitation
  \item Surgical intervention when indicated
  \item Regular monitoring and follow-up care
\end{itemize}
\section{Living with It}
\begin{itemize}[noitemsep]
  \item Follow your treatment plan as prescribed
  \item Maintain regular follow-up with your healthcare provider
  \item Make necessary lifestyle adjustments
  \item Seek support from family, friends, or support groups
  \item Stay informed about your condition
\end{itemize}
\section{Outlook}
\begin{itemize}[noitemsep]
  \item The outlook varies from person to person
  \item Migraine is a chronic condition with episodic attacks that can be managed but not cured.
\end{itemize}
